\section{Conclusion}
\label{sec:conclusion}

We trained three local ML models to accelerate SCF convergence in
MADNESS and all three failed. A level-clamped density predictor achieved
MSE ratios of 0.41--0.98$\times$ baseline at fine tree levels (10--14)
and 1.00--1.12$\times$ at coarse levels (1--9), producing identical SCF
iteration counts across all 14 molecule-threshold combinations.
A parent-node feature model expanded the architecture to 4.1M parameters
and doubled the input dimension, yet reproduced the same per-level MSE
ratios as the base model, confirming that additional capacity extracts no
new signal from the available inputs. A tree structure predictor reached
F1\,=\,0.860 but required 31 SCF iterations versus 12 for the
promolecular guess, a 2.6$\times$ degradation.

The common root cause is an input signal bottleneck at coarse tree
levels. MADNESS converges its SCF solution at levels 1--9 before
refining finer scales, so only coarse-level density corrections can
reduce iteration counts. The promolecular density and nuclear potential
at those levels encode atomic superposition and electrostatics but carry
no electron-electron interaction information. No local model operating on
these inputs can produce the exchange-correlation corrections needed to
improve the coarse-level density, regardless of architecture or training
procedure.

This mechanistic explanation (that multi-scale solvers gate convergence
at their coarsest resolution, where atom-centered inputs lack the
relevant physics) is not specific to MADNESS or to density functional
theory. It applies whenever an ML model receives only single-body inputs
and targets a quantity governed by many-body interactions, and the
iterative solver queries coarse scales first. Recognizing this pattern
before training avoids the cycle of architecture search and
hyperparameter tuning that preceded our diagnosis.

One path remains open. A converged Gaussian-basis density already
contains exchange-correlation physics and, when used directly as a
MADNESS starting guess, reduces SCF iteration counts. An ML model
trained on the small residual between the Gaussian-basis density and the
converged MRA density would operate on an input that does not suffer from
the signal bottleneck identified here. Whether this residual is
learnable with local features, and whether the resulting speedup
justifies the cost of the Gaussian-basis calculation, are questions that
require new experiments.
